\documentclass[11pt]{article}

\usepackage[margin=0.85in]{geometry}
\usepackage{amsmath,amssymb}
\usepackage{enumitem}
\usepackage{xcolor}
\usepackage{tcolorbox}
\usepackage{fancyhdr}
\usepackage{lastpage}
\usepackage{hyperref}
\usepackage{microtype}

\definecolor{uwpurple}{HTML}{4B2E83}
\definecolor{uwgold}{HTML}{B7A57A}
\definecolor{softpurple}{HTML}{F4F1FA}

\hypersetup{colorlinks=true, linkcolor=uwpurple, urlcolor=uwpurple}

\pagestyle{fancy}
\fancyhf{}
\lhead{\textbf{MATH 394: Probability I}}
\rhead{\textbf{Homework 7}}
\cfoot{\thepage\ of \pageref{LastPage}}

\setlength{\parindent}{0pt}
\setlength{\parskip}{0.45em}

\newcommand{\course}{MATH 394: Probability I}
\newcommand{\instructor}{Arman Jahangiri}
\newcommand{\department}{Department of Mathematics}
\newcommand{\university}{University of Washington}

\newtcolorbox{problemheader}{
    colback=softpurple,
    colframe=uwpurple,
    boxrule=0.8pt,
    arc=2mm,
    left=2mm,
    right=2mm,
    top=1mm,
    bottom=1mm
}

\newcommand{\source}[2]{%
    \vspace{0.15cm}
    {\small\textit{Source: Blitzstein and Hwang, \textit{Introduction to Probability}, Chapter #1, Problem #2.}}
}

\newcommand{\answerbox}{%
\begin{flushright}
\begin{tcolorbox}[width=0.36\textwidth,height=2cm,colback=white,colframe=uwpurple,boxrule=0.7pt,arc=1mm]
{\small\textbf{Final Answer}}
\end{tcolorbox}
\end{flushright}
}

\newcommand{\partanswerbox}[1]{%
\begin{flushright}
\begin{tcolorbox}[width=0.36\textwidth,height=2cm,colback=white,colframe=uwpurple,boxrule=0.7pt,arc=1mm]
{\small\textbf{Final Answer for Part #1}}
\end{tcolorbox}
\end{flushright}
}

\begin{document}

\begin{titlepage}
\centering
\vspace*{1.2cm}

{\Large \textbf{\university}}\\
{\large \department}

\vspace{1.5cm}

{\Huge \color{uwpurple}\textbf{Homework 7}}\\[0.25cm]
{\Large \textbf{\course}}\\[0.25cm]
{\large Instructor: Arman Jahangiri}

\vspace{1cm}

\begin{tcolorbox}[width=0.78\textwidth,colback=softpurple,colframe=uwpurple,boxrule=0.9pt,arc=3mm]
\centering
\textbf{Submission:} Single PDF on Gradescope

\textbf{Deadline:} 10:00 PM Pacific Time on the listed due date
\end{tcolorbox}

\vspace{0.8cm}

\begin{tcolorbox}[width=0.86\textwidth,colback=white,colframe=uwgold,boxrule=0.8pt,arc=3mm]
This homework is worth \textbf{50 points}.

Across the quarter, there are \textbf{eight homework assignments}, worth \textbf{400 points total}. Homework assignments together account for \textbf{40\%} of the final course grade. Further course policies can be seen in the \href{https://armanjg.github.io/MATH-394-Probability-1-v2/}{\textbf{MATH 394 syllabus}.}
\end{tcolorbox}

\vfill
\end{titlepage}

\section*{Homework Policy}

\subsection*{Submission and Deadline}

Please:

\begin{itemize}[leftmargin=*]
    \item Submit homework as a single PDF on Gradescope by \textbf{10:00 PM (Pacific Time)} on the listed due date.
    \item Begin each problem on a new page, and ensure that pages are correctly tagged in Gradescope.
    \item Write solutions clearly and legibly.
    \item Typesetting solutions in LaTeX is encouraged, but not required.
\end{itemize}

\subsection*{Late Days}

Each student is allotted \textbf{six late days} for the quarter. A late day extends a homework deadline by up to 24 hours without penalty. For example, submitting an assignment anytime between 10:01 PM on the due date and 10:00 PM the following day counts as one late day.

The following rules apply:
\begin{itemize}[leftmargin=*]
   \item No explanation is required to use late days.
   \item At most \textbf{two late days} may be used on any single homework assignment.
   \item Late days may not be used on the final homework assignment or on any assignment for which solutions have already been released.
   \item It is the student's responsibility to keep track of how many late days have been used during the quarter.
\end{itemize}

Once all late days have been exhausted, additional late submissions will incur a penalty of \textbf{10\% per day}, up to a maximum deduction of \textbf{50\%}. Assignments submitted more than five days late, or after solutions have been released, will not be accepted.

\subsection*{Submission Issues and Technical Difficulties}

If a serious technical issue prevents a timely Gradescope submission, students may temporarily submit their assignment by email to the instructor at \href{mailto:armanjg@uw.edu}{armanjg@uw.edu}. In such cases:

\begin{itemize}[leftmargin=*]
    \item The submission must consist of a \textbf{single PDF file}. The email subject line must follow the format:
    \[
    \texttt{HW\# Submission - FULL NAME - STUDENT ID}.
    \]
    \item Once the Gradescope issue has been resolved, the student must upload the \emph{exact same file} to Gradescope as soon as possible.
    \item The student should clearly indicate in the Gradescope submission that the assignment was previously submitted by email before the deadline.
\end{itemize}

Email submissions are intended only for genuine technical emergencies and should not be used as a substitute for timely Gradescope submission.

\newpage

% ------------------------------------------------------------
\begin{problemheader}
\textbf{Problem 1.  }
\end{problemheader}
\source{7}{1}

Alice and Bob arrange to meet for lunch on a certain day at noon. However, neither is known for punctuality. They both arrive independently at uniformly distributed times between noon and 1 pm on that day. Each is willing to wait up to 15 minutes for the other to show up.

What is the probability they will meet for lunch that day?

\newpage

% ------------------------------------------------------------

\begin{problemheader}
\textbf{Problem 2.  }
\end{problemheader}
\source{7}{7}

A stick of length $L$, where $L$ is a positive constant, is broken at a uniformly random point $X$. Given that $X=x$, another breakpoint $Y$ is chosen uniformly on the interval $[0,x]$.

\begin{enumerate}[label=(\alph*), leftmargin=*]
    \item Find the joint PDF of $X$ and $Y$. Be sure to specify the support.

    \item We already know that the marginal distribution of $X$ is $\operatorname{Unif}(0,L)$. Check that marginalizing out $Y$ from the joint PDF agrees that this is the marginal distribution of $X$.

    \item We already know that the conditional distribution of $Y$ given $X=x$ is $\operatorname{Unif}(0,x)$. Check that using the definition of conditional PDFs agrees that this is the conditional distribution of $Y$ given $X=x$.

    \item Find the marginal PDF of $Y$.

    \item Find the conditional PDF of $X$ given $Y=y$.
\end{enumerate}

\newpage

% ------------------------------------------------------------
 
\begin{problemheader}
\textbf{Problem 3.}
\end{problemheader}
\source{7}{15}

Let $X$ and $Y$ be continuous random variables, with joint CDF $F(x,y)$. Show that the probability that $(X,Y)$ falls into the rectangle $[a_1,a_2]\times[b_1,b_2]$ is
\[
F(a_2,b_2)-F(a_1,b_2)+F(a_1,b_1)-F(a_2,b_1).
\]


\newpage

% ------------------------------------------------------------

\begin{problemheader}
\textbf{Problem 4.  }
\end{problemheader}
\source{7}{16}

Let $X$ and $Y$ have joint PDF
\[
f_{X,Y}(x,y)=x+y,\qquad 0<x<1,\;0<y<1.
\]

\begin{enumerate}[label=(\alph*), leftmargin=*]
    \item Check that this is a valid joint PDF.

    \item Are $X$ and $Y$ independent?

    \item Find the marginal PDFs of $X$ and $Y$.

    \item Find the conditional PDF of $Y$ given $X=x$.
\end{enumerate}

\newpage

% ------------------------------------------------------------

\begin{problemheader}
\textbf{Problem 5. }
\end{problemheader}
\source{7}{17}

Let $X$ and $Y$ have joint PDF
\[
f_{X,Y}(x,y)=cxy,\qquad 0<x<y<1.
\]

\begin{enumerate}[label=(\alph*), leftmargin=*]
    \item Find $c$ to make this a valid joint PDF.

    \item Are $X$ and $Y$ independent?

    \item Find the marginal PDFs of $X$ and $Y$.

    \item Find the conditional PDF of $Y$ given $X=x$.
\end{enumerate}

\newpage

% ------------------------------------------------------------
\begin{problemheader}
\textbf{Problem 6.}
\end{problemheader}
\source{7}{21}

Find the probability that the quadratic polynomial
$$
Ax^2+Bx+1,
$$
where the coefficients $A$ and $B$ are determined by drawing i.i.d. $\operatorname{Unif}(0,1)$ random variables, has at least one real root.

\emph{Hint: By the quadratic formula, the polynomial $Ax^2+Bx+C$ has a real root if and only if $B^2-4AC\ge0$. Try to work on the event $\{B^2-4AC\geq 0\}\in \mathcal{F}$. } 

\newpage

% ------------------------------------------------------------


\begin{problemheader}
\textbf{Problem 7. }
\end{problemheader}
\source{7}{39--40}

\begin{enumerate}[label=(\alph*), leftmargin=*]
    \item Two fair, six-sided dice are rolled, one green and one orange, with outcomes $X$ and $Y$ for the green die and the orange die, respectively.
    \begin{enumerate}[label=(\roman*), leftmargin=*]
        \item Compute the covariance of $X+Y$ and $X-Y$.
        \item Are $X+Y$ and $X-Y$ independent?
    \end{enumerate}

    \item Let $X$ and $Y$ be i.i.d. $\operatorname{Unif}(0,1)$.
    \begin{enumerate}[label=(\roman*), leftmargin=*]
        \item Compute the covariance of $X+Y$ and $X-Y$.
        \item Are $X+Y$ and $X-Y$ independent?
    \end{enumerate}
\end{enumerate}

\newpage

% ------------------------------------------------------------

\begin{problemheader}
\textbf{Problem 8. }
\end{problemheader}
\source{7}{65}

Let
\[
(X_1,\dots,X_k)
\]
be Multinomial with parameters
\[
n \quad \text{and} \quad (p_1,\dots,p_k).
\]
Use indicator random variables to show that
\[
\operatorname{Cov}(X_i,X_j)=-np_ip_j,
\qquad i\ne j.
\]
\end{document}